%% ch08-scaling.tex — JA4proxy Operator's User Guide
%% Chapter 8: Scaling

\chapter{Scaling}
\label{ch:scaling}
\index{scaling}

\newthought{Scaling \ja\ is straightforward because the proxy is stateless per-instance.}
All shared state (bans, blacklists, rate-limit windows, enrichment caches) lives in
Redis, which every proxy instance reads and writes. Adding more proxy instances means
adding more connection processing capacity without any coordination overhead between
instances.

This chapter covers horizontal scaling with Docker Compose, the Kubernetes deployment
option, and guidance on when to consider the Go proxy.

\section{Performance Baseline}
\index{performance!baseline}

Before scaling, it helps to understand the single-instance baseline:

\begin{table}[ht]
  \small
  \begin{tabularx}{\linewidth}{lX}
    \toprule
    \textbf{Configuration} & \textbf{Throughput} \\
    \midrule
    Single Go instance, bridged container networking & $\sim$600 connections/second \\
    Single Go instance, host networking & $\sim$3{,}000 connections/second \\
    Horizontal scale-out across $N$ nodes & scales $\sim$linearly with $N$ \\
    Python prototype (historical, not for production) & $\sim$350 connections/second \\
    \bottomrule
  \end{tabularx}
  \caption{Measured end-to-end throughput, full scoring pipeline, on commodity
           hardware (Intel i9-9900K). See \code{docs/performance/benchmarks.md}.}
\end{table}

The Go engine itself is not the bottleneck. Under bridged container networking the
limit is Docker's userland \code{docker-proxy} relay ($\sim$600 conn/s); deploying
with host networking removes it and the same code sustains $\sim$3{,}000 conn/s.
A single modest server therefore handles several hundred connections per second
comfortably; scale horizontally beyond that.

\section{Docker Compose Scaling}
\index{scaling!Docker Compose}
\index{scale-proxies.sh}

\subsection{Scaling to Multiple Instances}

\begin{bashcode}
./scripts/scale-proxies.sh 4
\end{bashcode}

This adjusts the Docker Compose configuration to run 4 proxy instances and
reloads HAProxy's backend configuration to load-balance across them. You should
see all 4 instances appear in the HAProxy statistics page at
\url{http://localhost:8404/stats}.

HAProxy uses round-robin load balancing across proxy instances. Each proxy
instance registers its own metrics endpoint; Prometheus is pre-configured
to scrape all of them via the Docker service discovery mechanism.

\subsection{What Is and Is Not Shared}

\begin{table}[ht]
  \small
  \begin{tabularx}{\linewidth}{lXl}
    \toprule
    \textbf{State} & \textbf{Where stored} & \textbf{Shared across instances?} \\
    \midrule
    IP bans & Redis (ban:* keys) & Yes — all instances see all bans \\
    JA4 blacklist & Redis SET + in-process & Yes — synced from Redis at startup and reload \\
    JA4 whitelist & Redis SET + in-process & Yes \\
    Rate-limit windows & Redis Sorted Sets & Yes — truly cluster-wide sliding windows \\
    Beaconing history & Redis Sorted Sets & Yes \\
    GeoIP database & Read-only mmap file & No — each instance reads its own copy \\
    Local decision cache & In-process LRU & No — per-instance, provides fast-path for repeat visitors \\
    Enrichment cache & Redis & Yes \\
    Analytics findings & Redis & Yes \\
    \bottomrule
  \end{tabularx}
  \caption{State sharing between proxy instances}
\end{table}

The local in-process cache is per-instance and not shared. This means a client
whose connection was allow-cached on instance A may be re-scored on instance B.
This is acceptable: the cache is an optimisation, not a correctness requirement.
The first-time score is the authoritative result; the cache prevents re-scoring
the same fingerprint+IP combination on the same instance.

\subsection{HAProxy Configuration}

HAProxy's configuration is in \code{config/haproxy.cfg}. For the development
stack, it is auto-generated by \code{scripts/scale-proxies.sh}. For production
deployments, you will manage this directly.

The key section for proxy backend configuration:

\begin{lstlisting}[language=bash, numbers=none]
backend ja4proxy_backend
    balance roundrobin
    option tcp-check
    server ja4proxy-1 ja4proxy-1:8080 check inter 2s fall 3 rise 2
    server ja4proxy-2 ja4proxy-2:8080 check inter 2s fall 3 rise 2
    server ja4proxy-3 ja4proxy-3:8080 check inter 2s fall 3 rise 2
    server ja4proxy-4 ja4proxy-4:8080 check inter 2s fall 3 rise 2
\end{lstlisting}

HAProxy performs TCP health checks every 2 seconds. If 3 consecutive checks fail,
the instance is removed from rotation. When 2 consecutive checks pass, it is
returned to rotation. This provides fast failure detection without flapping.

\section{Redis Scaling Considerations}
\index{Redis!scaling}

A single Redis instance handles the shared state for all proxy instances. At 4
proxy instances with Redis pipeline batching, Redis is typically below 10\% CPU
utilisation. At 8+ instances or very high connection rates, consider:

\begin{itemize}
  \item Enabling the Unix domain socket connection (reduces per-operation overhead)
  \item Increasing Redis's \code{maxmemory} and enabling an appropriate eviction
    policy for enrichment caches
  \item For production deployments: Redis Sentinel for HA, or Redis Cluster for
    horizontal scaling
\end{itemize}

\begin{warningbox}[Redis is a single point of failure]
In the development stack, a single Redis instance handles all shared state.
If Redis becomes unavailable, the proxy fails open (allows all connections)
rather than failing closed (blocking everything). This is the correct behaviour
for a security proxy, but it means your defence is weakened during Redis downtime.

For production: deploy Redis with Sentinel (3 nodes) for automatic failover.
\end{warningbox}

\section{Kubernetes Deployment}
\index{Kubernetes}
\index{Helm}

\ja\ ships with a Helm chart in \code{deploy/helm/ja4proxy/}.

\subsection{Basic Installation}

\begin{bashcode}
helm install ja4proxy ./deploy/helm/ja4proxy \
  --set proxy.backendHost=prod.internal.example.com \
  --namespace ja4proxy \
  --create-namespace
\end{bashcode}

\subsection{Key Helm Values}

\begin{table}[ht]
  \small
  \begin{tabularx}{\linewidth}{llX}
    \toprule
    \textbf{Value} & \textbf{Default} & \textbf{Description} \\
    \midrule
    \code{proxy.backendHost} & (required) & Hostname of backend server \\
    \code{proxy.dialSetting} & \code{0} & Initial dial value (0 = monitor) \\
    \code{proxy.replicas} & \code{2} & Initial replica count \\
    \code{hpa.enabled} & \code{true} & Enable horizontal pod autoscaler \\
    \code{hpa.minReplicas} & \code{2} & Minimum pod count \\
    \code{hpa.maxReplicas} & \code{20} & Maximum pod count \\
    \code{hpa.targetCPU} & \code{70} & CPU utilisation target (\%) \\
    \code{redis.external} & \code{true} & Use external Redis (recommended) \\
    \code{redis.host} & \code{redis} & Redis hostname \\
    \code{redis.port} & \code{6379} & Redis port \\
    \code{resources.requests.cpu} & \code{500m} & CPU request per pod \\
    \code{resources.requests.memory} & \code{256Mi} & Memory request per pod \\
    \code{resources.limits.cpu} & \code{2000m} & CPU limit per pod \\
    \code{resources.limits.memory} & \code{512Mi} & Memory limit per pod \\
    \bottomrule
  \end{tabularx}
  \caption{Key Helm chart values}
\end{table}

\subsection{Production Kubernetes Setup}

For a production Kubernetes deployment:

\begin{bashcode}
helm install ja4proxy ./deploy/helm/ja4proxy \
  --set proxy.backendHost=prod.internal.example.com \
  --set proxy.dialSetting=0 \
  --set proxy.replicas=3 \
  --set hpa.maxReplicas=20 \
  --set redis.external=true \
  --set redis.host=redis.prod-infra.svc.cluster.local \
  --set resources.requests.cpu=500m \
  --set resources.limits.cpu=2000m \
  --namespace ja4proxy \
  --create-namespace
\end{bashcode}

\begin{notebox}[External Redis for Kubernetes]
Set \configkey{redis.external=true} and provide your own Redis instance (Redis
Sentinel or Redis Cluster) for Kubernetes deployments. The bundled Redis in the
Helm chart is for development only — it has no persistence or HA configuration.
\end{notebox}

\subsection{Horizontal Pod Autoscaler}

The HPA is configured to scale on CPU utilisation. At 70\% CPU across the pod set,
Kubernetes adds pods up to \configkey{hpa.maxReplicas}. The scale-down cooldown is
5 minutes to prevent thrashing.

Monitor HPA status:
\begin{bashcode}
kubectl get hpa -n ja4proxy
kubectl describe hpa ja4proxy -n ja4proxy
\end{bashcode}

\subsection{Configuration in Kubernetes}

The \code{config/proxy.yml} content is mounted as a ConfigMap. To update
configuration in Kubernetes:

\begin{bashcode}
# Update the ConfigMap
kubectl create configmap ja4proxy-config \
  --from-file=proxy.yml=config/proxy.yml \
  --namespace ja4proxy \
  --dry-run=client -o yaml | kubectl apply -f -

# Trigger a rolling reload by annotating the deployment
kubectl rollout restart deployment/ja4proxy -n ja4proxy
\end{bashcode}

\begin{notebox}[Hot reload in Kubernetes]
Hot reload via SIGHUP works in Kubernetes if you exec into the pod. For a cluster-wide
config reload (all pods), it is simpler to do a rolling restart as shown above. A
rolling restart replaces pods one at a time, maintaining availability.
\end{notebox}

\section{Scaling Horizontally}
\index{scaling!horizontal}

The Go engine is the production runtime, and instances are stateless — the only
per-instance state is the in-process LRU cache, which is an optimisation, not a
source of truth. Scaling out is therefore just running more instances behind your
load balancer; bans and findings are shared through Redis, so every instance sees
the same view.

\begin{itemize}
  \item Add instances when sustained load approaches a single node's ceiling
    (several hundred conn/s under bridged networking, $\sim$3{,}000 with host
    networking — see the baseline table above).
  \item Point your load balancer (or the outside firewall directly) at the
    instances; no shared local state needs synchronising.
  \item Size Redis for the aggregate signal-write rate across all instances.
\end{itemize}
